\documentclass{article}
\usepackage[T1]{fontenc}
\usepackage{lmodern}
\usepackage{graphicx}
\usepackage{amsmath,amssymb}
\usepackage[a4paper,margin=2cm]{geometry}
\usepackage[absolute,overlay]{textpos}
\setlength{\TPHorizModule}{1mm}
\setlength{\TPVertModule}{1mm}
\usepackage[indonesian]{babel}
\usepackage{hyperref}
\setlength{\emergencystretch}{2em}
% unit: Halaman 102

\begin{document}

% === Halaman 102 (python/native) ===

• Untuk matriks 3 x 3, matriks dapat dicari dengan metode Eliminasi Gauss-Jordan, 

atau dihitung langsung dengan rumus berikut:: 

• yang dalam hal ini,

    A = (ei – hf)         B = – (di – fg)     C = (dh – eg) 

    D = – (bi – hc)     E = (ai – cg)        F = – (ah – bg)

    G = (bf – ec) H = – (af – cd)           I = (ae – bd) 

dan


det(K) = aA + bB + cC

102

K = 





















i

h

g

f

e

d

c

b

a

,   K–1 = 





















=





















I

F

C

H

E

B

G

D

A

I

H

G

F

E

D

C

B

A

T

K)

K)

det(

1

det(

1

\end{document}
